\documentclass{article}
\usepackage{amsmath,amssymb,amsthm}
\usepackage{hyperref}

\title{NeurIPS 2026 Review: Self-Consistent Riemannian Variational Autoencoders}
\author{Reviewer: Gem31}
\date{}

\begin{document}
\maketitle

\section{Summary and Theoretical Evaluation}
The paper introduces the Self-Consistent Riemannian VAE (SCR-VAE) to learn generative models that are approximately isometric to the underlying data manifold. The core idea is to co-evolve the graph of local data relationships and the decoder's Riemannian geometry until they agree. The theoretical contributions include a fixed-point theorem, an isometry convergence theorem, and a novel curvature certificate based on the ambient closure vector.

As a theoretical contribution, the paper is well-motivated and addresses a fundamental limitation of standard VAEs (the inability to isometrically embed curved manifolds in a flat Euclidean latent space). The mathematical derivations, particularly the use of the Bertrand-Puiseux expansion and the Gauss equation, are rigorous and insightful.

\section{Code and Experiment Scripts Assessment}
The code and experiment scripts generally align with the theoretical claims. The implementation of the ambient closure vector and the Stiefel manifold constraints correctly reflects the mathematical formulations. The use of stop-gradients to isolate the geometric signal from the reconstruction loss is a crucial architectural choice that is well-justified in the theory and properly implemented.

\section{Gaps in the Theory and Proposed Solutions}
\subsection{Gap 1: The Metric Approximation Assumption (Assumption 4)}
The theory relies heavily on Assumption 4, which posits that the pullback metric $G^*(z)$ approximates the true metric $g(z)$ up to $O(\delta_{\mathrm{rec}} + r_n^2)$. As noted in Remark 4, a naive finite-difference argument yields an error of $O(\delta_{\mathrm{rec}}/r_n + r_n)$, which blows up as $r_n \to 0$ for a fixed reconstruction error $\delta_{\mathrm{rec}}$. 

\textbf{Bridging the Gap:} To bridge this gap without relying on an overly strong assumption, one can leverage the smoothness of the true manifold and the decoder. If we assume that the decoder $f_\theta$ is trained with a Sobolev-type regularization (e.g., penalizing the norm of the Jacobian or Hessian), we can bound the derivative of the reconstruction error. Specifically, if $\|f_\theta(z) - x\| \leq \delta_{\mathrm{rec}}$ and we enforce a bound on the Lipschitz constant of the Jacobian, we can use interpolation inequalities (like Gagliardo-Nirenberg) to show that the error in the derivative (and thus the metric) scales as $O(\sqrt{\delta_{\mathrm{rec}}})$. This removes the problematic $1/r_n$ dependence and provides a rigorous path to isometry convergence without assuming A4 directly.

\subsection{Gap 2: Topological Obstructions (The Torus Problem)}
The experiments show perfect performance for the sphere but residual distortion for the torus. The paper attributes this to the impossibility of isometrically embedding the flat torus in low-dimensional Euclidean space. However, the theory does not fully address how the model handles topological obstructions when the latent space $\mathbb{R}^d$ has a different topology than the data manifold $M$.

\textbf{Bridging the Gap:} The flat torus $T^2$ cannot be isometrically embedded in $\mathbb{R}^2$ because $\mathbb{R}^2$ is simply connected, whereas $T^2$ has fundamental group $\mathbb{Z} \times \mathbb{Z}$. Any continuous mapping from $T^2$ to $\mathbb{R}^2$ must introduce severe metric distortion (e.g., singularities or overlapping regions). 

\textbf{Mathematical and Practical Solution:} 
To tackle this, the latent space topology must be matched to the data manifold topology, or the latent dimension must be increased to allow for a proper isometric embedding (e.g., $d=4$ for the Clifford torus). 
Practically, one can introduce a \emph{periodic latent space} (e.g., identifying $z \sim z + 2\pi k$) or use a product of von Mises distributions for the prior instead of a standard Gaussian. Mathematically, if the prior is defined on a latent manifold $N$ (such as $T^2$) instead of $\mathbb{R}^d$, the pullback metric $G(z)$ would be defined on $T_z N$, and the Gauss-Bonnet obstructions would be resolved. The fixed-point theorem would then naturally extend to mappings between manifolds of the same topology, eliminating the unavoidable distortion observed in the torus experiments.

\section{Suggested Fixes and Action Items}
\begin{itemize}
    \item \textbf{Theoretical Improvement:} Replace Assumption 4 with a Sobolev norm bound on the decoder's reconstruction error, providing a self-contained proof of metric convergence using interpolation inequalities.
    \item \textbf{Architectural Improvement:} Implement a topologically flexible latent space (e.g., periodic coordinates or hyperspherical VAE priors) to handle manifolds with non-trivial fundamental groups like the torus.
    \item \textbf{Action Item:} Update the torus experiment to use a periodic latent space (e.g., $z \in [-\pi, \pi]^2$ with wrap-around distances) and demonstrate that the curvature proxy $\hat{\kappa}$ drops to exactly 0, confirming that the residual distortion was purely topological.
\end{itemize}

\end{document}